% Generated draft: requires mathematical and source review.
\documentclass[11pt]{article}
\usepackage[margin=25mm]{geometry}
\usepackage{fontspec}
\setmainfont{Latin Modern Roman}
\usepackage{amsmath,amssymb,mathtools,unicode-math}
\setmathfont{Latin Modern Math}
\usepackage{enumitem,needspace,xurl,hyperref}
\setlength{\parindent}{0pt}
\setlength{\parskip}{5pt}
\setlength{\emergencystretch}{3em}
\newcommand{\N}{\mathbb{N}}
\newcommand{\Z}{\mathbb{Z}}
\newcommand{\Q}{\mathbb{Q}}
\newcommand{\R}{\mathbb{R}}
\newcommand{\C}{\mathbb{C}}
\newcommand{\F}{\mathbb{F}}
\newcommand{\A}{\mathbb{A}}
\newcommand{\PP}{\mathbb{P}}
\newcommand{\E}{\mathbb{E}}
\newcommand{\eps}{\varepsilon}
\newcommand{\dd}{\,\mathrm{d}}
\begin{document}
\section*{UnsolvedMath ID: 30003157}
\textbf{Source problem code:} catalogue entry
\subsection{Definitions and mathematical statement}
Let \(X\subset\PP^n\) be a hypersurface of degree \(d\). An integral curve is a reduced, irreducible curve; it is rational if it is birational to the projective line over the ground field. Its degree is measured with respect to \(\mathcal O_{\PP^n}(1)\). Write \(R_e(X)\) for the source's space of integral rational curves of degree \(e\) contained in \(X\). The source does not specify the ground field or the precise parameter-space construction underlying \(R_e(X)\).

The phrase ``very general hypersurface'' conventionally means a hypersurface whose parameter lies outside a countable union of proper closed subsets of the parameter space of degree-\(d\) hypersurfaces. The ground-field assumptions needed for this interpretation are not supplied.

\textbf{Conjectural statement as supplied.} For a very general degree-\(d\) hypersurface \(X\subset\PP^n\), the proposed equality is
\[
\dim R_e(X)=e(n-d+1)+n-4,
\qquad (n,d)\ne(3,4).
\]
The assertion is presented as applying throughout its intended parameter range, but that range is not specified. In particular, the source gives no bounds on \(n,d,e\), no nonemptiness hypothesis, and no convention for interpreting the equality when the right-hand side is negative. It also does not specify whether only the dimension of the entire space or the dimension of every irreducible component is intended. No replacement assertion is supplied for \((n,d)=(3,4)\). These unresolved qualifications prevent treating the displayed formula as a fully specified universal conjecture.
\subsection{Short English statement}
The supplied conjecture predicts that the space of degree-\(e\) integral rational curves on a very general degree-\(d\) hypersurface in \(\PP^n\) has dimension \(e(n-d+1)+n-4\), excluding \((n,d)=(3,4)\). The ground field, parameter ranges, parameter-space construction, and treatment of empty spaces remain unspecified. This is a statement-only draft requiring human mathematical review.
\subsection{Sources}
\begin{itemize}
\item Supplied local UnsolvedMath registry record, permanent ID \(30003157\), source problem code \texttt{catalogue entry}, title ``Expected Dimensions of Rational-Curve Spaces on Hypersurfaces,'' in the collection ``Oberwolfach Reports Open Problems.'' No individual proposer is identified; the supplied proposed year is 2016.
\item Supplied report citation: \emph{Classical Algebraic Geometry (2017)}, Oberwolfach Reports, pp.~1635--1682, DOI \texttt{10.4171/owr/2016/29}. \url{https://doi.org/10.4171/owr/2016/29}. This is also the supplied literature-source link labeled ``Oberwolfach Report 29.'' The background describes these reports as records of mathematical workshops and their open-problem discussions.
\item Supplied literature source: Riedl--Yang, \emph{Kontsevich's spaces of rational curves}. \url{https://arxiv.org/abs/1605.01759}.
\item Literature assessment dated exactly \texttt{2026-08-21}: the supplied curation labels the entry open and describes it as a broad collection of questions about rational curves, with results for important special families and general questions remaining active. This is supplied dated curation, not an independently verified literature assessment; the background expressly requests independent verification of time-sensitive status and cited claims.
\item The registry labels the statement \texttt{corrected\_verified} and reports manual semantic, typographic, and LaTeX review asserting self-containment and preservation of the mathematical target. Those are supplied editorial claims, not independent verification in this draft.
\end{itemize}
\end{document}
